\section{Related Work}
\label{sec:related}

\paragraph{Uplift and CATE benchmarks.} The dataset we use is the v2.1
release of the Criteo benchmark of Diemert et
al.~\cite{diemert2021}, who advise modeling uplift on visits because
conversions are too sparse, and who note that uplift models need very large
samples. We test both claims directly.

\paragraph{Meta-learners and forests.} The S-, T-, and X-learners are
surveyed and compared in K{\"u}nzel et al.~\cite{kunzel2019}. The DR-learner
follows Kennedy~\cite{kennedy2023}, and the transformed-outcome approach
follows Athey and Imbens~\cite{athey2015}. Honest causal forests are due to
Wager and Athey~\cite{wager2018}; we use the orthogonalized variant in
EconML~\cite{econml}, built on the double machine-learning framework of
Chernozhukov et al.~\cite{chernozhukov2018}.

\paragraph{Average-effect estimation.} For ATE we use regression adjustment
with full interactions and robust standard errors~\cite{lin2013}, inverse
propensity weighting~\cite{rosenbaum1983}, and augmented inverse-propensity
weighting (AIPW), which is doubly robust~\cite{robins1994}. The effect of
actually seeing an ad uses randomized assignment as an
instrument~\cite{imbens1994}.

\paragraph{Evaluating uplift.} Because the individual effect is never
observed, uplift rankings are scored with Qini and AUUC
curves~\cite{radcliffe2007}. We complement them with the best-linear-predictor
and sorted-group tests of Chernozhukov et al.~\cite{chernozhukov2018gates},
which ask whether estimated effects carry real heterogeneity and whether
they are calibrated.

\paragraph{Randomization checks.} To test whether treatment is predictable
from covariates we use a classifier two-sample test~\cite{lopezpaz2017}.

\paragraph{Advertising measurement.} Gordon et al.~\cite{gordon2019} compare
observational approaches to randomized field experiments in advertising and
find that adjustment often fails to recover the experimental answer. Our
observational demonstration (Section~\ref{sec:ate}) is in that spirit.

\paragraph{Gradient boosting.} LightGBM~\cite{ke2017} is our workhorse
base learner for every predictive and meta-learner model.
